\documentclass[11pt]{amsart}
\usepackage{amsmath,amssymb,amsthm,mathtools}
\usepackage[margin=1in]{geometry}
\usepackage{hyperref}
\usepackage{listings}
\lstset{basicstyle=\ttfamily\footnotesize,breaklines=true,frame=single,
  showstringspaces=false,columns=fullflexible,keepspaces=true}

\newtheorem{theorem}{Theorem}[section]
\newtheorem{proposition}[theorem]{Proposition}
\newtheorem{lemma}[theorem]{Lemma}
\newtheorem{corollary}[theorem]{Corollary}
\theoremstyle{definition}
\newtheorem{definition}[theorem]{Definition}
\theoremstyle{remark}
\newtheorem{remark}[theorem]{Remark}
\newtheorem{example}[theorem]{Example}

\newcommand{\RH}{\mathrm{RH}}

\title[Eventual versus uniform positivity]{Eventual versus uniform positivity in
scale-dependent hierarchies}
\author{David Alejandro Trejo Pizzo}
\thanks{Independent Researcher. \texttt{dtrejopizzo@gmail.com}}


\begin{document}
\begin{abstract}
We separate eventual and uniform positivity in scale-dependent positivity
hierarchies, record the elementary cases in which one implies the other, and
prove a continuity lemma, with an explicit rate, for the Hodge--Riemann gap of a
convergent family of forms. Several examples (discrete and mixed moment measures,
toy families of forms) show what convergence to a positive limit does and does
not give. We discuss heuristically how the Jensen-polynomial hierarchy of $\xi$
fits this language; no statement about the Riemann Hypothesis is proved. A
numerical appendix with code illustrates the examples.
\end{abstract}
\maketitle

\tableofcontents
\newpage

\section{Introduction}
Several reformulations of RH are positivity statements of infinite order:
the hyperbolicity of all Jensen polynomials $J_{d,n}$ of $\xi$ (P\'olya), and
the total positivity of the function $\Lambda$ with Laplace transform $1/\Xi$
\cite{RP06}. For the Jensen hierarchy the finite orders are partly known: each
fixed degree is hyperbolic beyond an effective shift~\cite{GORZ}, while
hyperbolicity at all degrees and shifts is equivalent to RH;
see~\cite{RP06,RP14}. For $\Lambda$ we know of no unconditional finite-order
result, and we do not use it further. This paper isolates the quantifier
structure involved (eventual versus uniform positivity), records the elementary
situations in which the two coincide, and proves a continuity lemma for the
Hodge--Riemann gap of a convergent family of forms. The connection with $\zeta$
is discussed only heuristically.

\section{Scale-parametrized positivity hierarchies}
\begin{definition}
A \emph{scale-parametrized positivity hierarchy} is a family of real numbers
$\delta_k(T)$ ($k\ge1$, $T\ge1$), nondecreasing in $k$:
$\delta_k(T)\le\delta_{k+1}(T)$. We write $P_k(T)$ for the statement
$\delta_k(T)\le0$; by monotonicity $P_{k+1}(T)$ implies $P_k(T)$.
\end{definition}
Two statements must be separated:
\[
\textbf{(eventual)}\quad\forall k\ \exists T_0(k)\ \forall T>T_0(k):P_k(T);
\qquad
\textbf{(uniform)}\quad\forall T\ \forall k:P_k(T).
\]
When \textbf{(eventual)} holds we let $T_0(k)\ge1$ denote the infimum of the
admissible thresholds in $[1,\infty)$. \textbf{(uniform)} is the
statement $\sup_k\sup_T\delta_k(T)\le0$.

\begin{remark}
The finite-to-full gap is \emph{not} the implication $\forall k\,P_k(T)\Rightarrow
P_\infty(T)$, which is trivial when $P_\infty(T)$ is defined as the conjunction
of all $P_k(T)$. It is the quantifier swap: each $P_k$ may be known only eventually
in $T$, with no control of $T_0(k)$ uniform in $k$. For $\xi$ we use the Jensen hierarchy: writing
$\Xi(t)=\xi(\tfrac12+it)=\sum_{k\ge0}\gamma(k)(-t^2)^k/k!$, the Jensen
polynomials are $J_{d,n}(X)=\sum_{j=0}^d\binom dj\gamma(n+j)X^j$ (denoted
$J^{d,n}$ in \cite{RP14}); order $k$ is the hyperbolicity of the
degree-$k$ Jensen polynomials $J_{k,n}$, the scale is the shift $n$; degree $2$
is the Tur\'an inequalities, each degree $k$ is known beyond an effective shift
by Griffin--Ono--Rolen--Zagier~\cite{GORZ}, and all degrees and shifts together
are RH. Monotonicity in $k$ at fixed $n$ holds because the derivative of the
reversed polynomial $X^{k+1}J_{k+1,n}(1/X)$ is $k+1$ times the reversal of
$J_{k,n}$, and differentiation and reversal preserve real-rootedness. For $\xi$
these known thresholds are thresholds of proofs: no bound uniform in $k$ is known,
and under RH every threshold may be taken to be the least scale.
\end{remark}

\section{The uniformity criterion}
\begin{theorem}[Criterion]\label{thm:crit}
Let $\delta_k(T)$ be a scale-parametrized hierarchy.
\begin{enumerate}
\item \emph{Finite type:} if $\delta_k(T)=\delta_k$ is independent of $T$, then
\textbf{(eventual)} $\iff$ \textbf{(uniform)}.
\item \emph{Bounded threshold:} if \textbf{(eventual)} holds and
$T_0:=\sup_kT_0(k)<\infty$, then $P_k(T)$ holds for all $k$ and all $T>T_0$.
\end{enumerate}
\end{theorem}
\begin{proof}
(1) \textbf{(uniform)} trivially implies \textbf{(eventual)}. Conversely,
\textbf{(eventual)} applied at any single $T>T_0(k)$ gives $\delta_k\le0$, and
$\delta_k$ does not depend on $T$. (2) For $T>T_0\ge T_0(k)$ the definition of
$T_0(k)$ gives $P_k(T)$, for every $k$.
\end{proof}
Case (2) gives uniform positivity only beyond a single scale; it says nothing
about scales $T\le T_0$. For instance $\delta_k(T)=1$ for $T<2$ and $\delta_k(T)=-1$
for $T\ge2$ has $T_0(k)=2$ for all $k$, yet \textbf{(uniform)} fails.

\section{The dividing line}\label{sec:line}
\begin{center}
\begin{tabular}{lccc}
\hline
hierarchy & scale parameter? & known threshold & finite-to-full \\
\hline
matroid Hodge--Riemann~\cite{AHK} & none & --- & scale-free: case (1) \\
Lorentzian polynomials~\cite{BH} & none & --- & scale-free: case (1) \\
$\xi$ Jensen polynomials & shift $n$ & $N(k)$, not uniform & \textbf{open} \\
\hline
\end{tabular}
\end{center}
Finite combinatorial objects are scale-free (case~(1)), so their full positivity
is a statement about each object separately (Adiprasito--Huh--Katz~\cite{AHK}
for matroids). For $\xi$, each degree $k$ is known to be hyperbolic beyond an
effective shift $N(k)$~\cite{GORZ}; no bound on $N(k)$ uniform in $k$ is known,
and under RH one may take $N(k)=0$. The finite-to-full gap for $\zeta$ is
therefore a gap in knowledge, not a property of the hierarchy that the present
framework can detect.

\section{An analogy with finitization}
\begin{remark}[heuristic]\label{cor:fin}
By Theorem~\ref{thm:crit}(1), a hierarchy that is finite type has no
finite-to-full gap. For the $\zeta$-hierarchies one may think of a
finite-dimensional model of the zeros (in the spirit of the Frobenius action in
the function-field case, cf.~\cite{RP06}) as something that would make the
hierarchy scale-free. We do not define such a ``finitization'' precisely here,
and we prove no equivalence between its absence and the finite-to-full gap; the
remark only records the analogy.
\end{remark}

\section{The bounded-threshold class is non-empty}
\begin{proposition}[Existence]\label{prop:exist}
There exist scale-parametrized positivity hierarchies that are not finite type
yet have bounded threshold.
\end{proposition}
\begin{proof}
Take $\delta_k(T)=-1/T$ for all $k$. Then $P_k(T)$ holds for all $k,T$, so
$T_0(k)=1$ for all $k$; and $\delta_k$ depends on $T$, so the hierarchy is not
finite type.
\end{proof}
The statement is elementary. Convergence of measures does not by itself produce
bounded thresholds in moment hierarchies: for any finite measure $\mu$
supported in $[-1,1]$, the smallest eigenvalue of the order-$N$ Hankel matrix
tends to $0$ as $N\to\infty$ (test the form on $x^N-x^{N-2}$, whose coefficient
vector has norm $\sqrt2$, while $\int x^{2N-4}(1-x^2)^2\,d\mu\to0$ by dominated
convergence). For Lebesgue measure on $[0,1]$ it is about $2.7\cdot10^{-3}$,
$1.1\cdot10^{-7}$, $1.1\cdot10^{-13}$ for matrix sizes $3,6,10$. So there is no
gap uniform in the order; only on a fixed order range $k\le K$ does a uniform gap,
and hence a perturbation argument, apply. See Example~\ref{ex:def}; Example~\ref{ex:deg} shows that even with bounded
threshold the gap need not be uniform.

\section{The convergence route}
One might hope that the convergence of the rescaled Jensen polynomials of $\xi$
to Hermite polynomials~\cite{GORZ} gives a passage to RH. We record only the
following.
\begin{remark}\label{thm:overRH}
If, in a fixed normalization, the coefficient vector of $J_{d,n}$ lies within the
real-rootedness radius of the rescaled Hermite polynomial for every $(d,n)$, then
every $J_{d,n}$ is hyperbolic and RH follows by the P\'olya--Jensen criterion.
Such a radius depends on the conditioning of the Hermite polynomial: the minimal
zero gap of $H_d$ is of order $d^{-1/2}$, not $1/d$ (numerically $0.96$, $0.69$,
$0.49$, $0.35$, $0.25$ for $d=5,10,20,40,80$), and the sensitivity of zeros to
coefficient perturbations is not controlled by the gap alone. Whether RH implies
such a bound is not addressed here.
\end{remark}

\section{Continuity of the Hodge--Riemann gap}
\begin{definition}[Hodge--Riemann tower]
A \emph{Hodge--Riemann tower} of degree $d$ is a sequence $(V_m,L_m,\deg_m)$, where
$V_m=\bigoplus_{j=0}^dV_m^j$ is a finite-dimensional graded commutative real
algebra with a fixed inner product, $L_m$ is multiplication by an element of
$V_m^1$,
$\deg_m\colon V_m^d\to\mathbb R$ is linear, and for $0\le j\le d/2$
\[
P_m^j=\ker\bigl(L_m^{\,d-2j+1}|_{V_m^j}\bigr),\qquad
Q_m^j(x,y)=\deg_m\bigl(x\cdot L_m^{\,d-2j}y\bigr)\quad(x,y\in V_m^j),
\]
with the Hodge--Riemann property $(-1)^jQ_m^j(v,v)>0$ for $0\ne v\in P_m^j$. The
\emph{HR gap} is $\delta_m=\min_j\inf_{v\in P_m^j,\,\|v\|=1}(-1)^jQ_m^j(v,v)$,
the minimum over those $j$ with $P_m^j\ne0$; since
each $P_m^j$ is finite-dimensional, $\delta_m>0$ for every $m$. The tower
\emph{converges} to $(V_\infty,L_\infty,Q_\infty^j)$ if all $V_m^j=V_\infty^j$ are
identified with a common space, $Q_m^j\to Q_\infty^j$ in operator norm, and the
orthogonal projections $\Pi_m^j$ onto $P_m^j$ converge in operator norm to those
$\Pi_\infty^j$ onto $P_\infty^j=\ker(L_\infty^{\,d-2j+1}|_{V_\infty^j})$; we
assume $P_\infty^j\ne0$ for the $j$ considered (then $P_m^j\ne0$ for large $m$,
since the projections converge), and $\delta_\infty$ is defined by the same
formula (it may be $0$). We write $\|Q_m-Q_\infty\|=\max_j\|Q_m^j-Q_\infty^j\|$.
\end{definition}

\begin{lemma}[Continuity of the HR gap]\label{thm:hr}
For a convergent Hodge--Riemann tower, $\delta_m\to\delta_\infty$. In particular
$\inf_m\delta_m>0$ if and only if $\delta_\infty>0$.
\end{lemma}
\begin{proof}
The convergence $\delta_m\to\delta_\infty$ follows from
Proposition~\ref{prop:rate} with
\[\eta(m)=\max\bigl(\|Q_m-Q_\infty\|,\ \max_j\|\Pi_m^j-\Pi_\infty^j\|\bigr)\to0.\] If $\delta_\infty>0$, then
$\delta_m\ge\delta_\infty/2$ for $m\ge m_0$, and the finitely many
$\delta_m$, $m<m_0$, are positive, so $\inf_m\delta_m>0$. If $\delta_\infty=0$,
then $\delta_m\to0$ and $\inf_m\delta_m=0$.
\end{proof}

\begin{corollary}\label{cor:towerbridge}
In a convergent Hodge--Riemann tower positivity holds at every level whether or
not $\delta_\infty>0$; only the uniformity of the gap depends on the limit.
\end{corollary}

\begin{remark}
Lemma~\ref{thm:hr} is a standard perturbation statement and uses no
combinatorial input. Adiprasito--Huh--Katz~\cite{AHK} prove the Hodge--Riemann
property for each matroid; the lemma says only how the gap behaves in a convergent
family. Example~\ref{ex:deg} shows that a degenerate limit is compatible with
positivity at every level.
\end{remark}

\section{Worked examples}
\begin{example}[Weak convergence without bounded threshold]\label{ex:def}
Let $\mu_m=\tfrac1m\sum_{i=1}^m\delta_{i/m}$ on $[0,1]$, converging weakly (not
in total variation) to Lebesgue measure, whose Hankel forms
$H^{(N)}=[\int x^{a+b}\,d\mu]_{0\le a,b\le N}$ are positive definite for every $N$.
Here the scale is $m$ and $P_N(m)$ is the statement that $H_m^{(N)}$ is
positive definite (defect $\delta_N(m)=-\lambda_{\min}(H_m^{(N)})$ if
$\lambda_{\min}(H_m^{(N)})>0$, and $\delta_N(m)=1$ otherwise; it is
nondecreasing in $N$ by eigenvalue interlacing). Since $\mu_m$ has $m$ atoms,
$H_m^{(N)}$ is positive definite if and only if $N<m$. Hence positivity at
order $N$ holds exactly for $m\ge N+1$, so $T_0(N)=N$, which is unbounded: weak
convergence to a measure with positive definite Hankel forms does not give a
bounded threshold. (Numerically, for $m=4$ the smallest eigenvalue is
$1.9\cdot10^{-5}$ at $N=3$ and $0$ up to rounding at $N=4$.) By contrast,
$\mu_m=(1+\tfrac1m)\lambda_{[0,1]}$ has a bounded threshold trivially.
\end{example}

\begin{example}[No uniform gap, positivity at every level]\label{ex:deg}
Let $\nu_m=(1-\tfrac1m)\delta_0+\tfrac1m\,\lambda_{[0,1]}$, converging to
$\delta_0$, whose Hankel forms have rank one. Since $\nu_m$ has a component with
positive density, every Hankel form of $\nu_m$ is positive definite for every
$m\ge2$, so the threshold is bounded. But at fixed order the smallest eigenvalue
is $\asymp1/m$ (proved below; illustrated in Appendix~\ref{app:deg}): the positivity gap is not
uniform in $m$, although positivity holds at every $m$. Indeed
$H^{(N)}(\nu_m)=(1-\tfrac1m)e_0e_0^{T}+\tfrac1mH^{(N)}(\lambda)$, so the smallest
eigenvalue lies between $\tfrac1m\lambda_{\min}(H^{(N)}(\lambda))$ and
$\tfrac1m\lambda_{\min}$ of the Hankel form of $\lambda$ restricted to $e_0^\perp$.
\end{example}

\begin{example}[Uniform matroids]\label{ex:matroid}
\cite{AHK} gives the Hodge--Riemann property for each uniform matroid $U_{r,m}$
individually. For $r\ge3$ the Chow rings of $U_{r,m}$ have dimension growing with
$m$, so they cannot be identified with one fixed finite-dimensional space and
Lemma~\ref{thm:hr} does not apply as stated; a uniform Hodge--Riemann statement
would need an infinite-dimensional version of the lemma, which we do not have.
\end{example}

\section{The \texorpdfstring{$\zeta$}{zeta} orders, concretely}
For $\xi$ the hierarchy is the hyperbolicity of the Jensen polynomials built
from the moments of the kernel $\Phi$
($\Xi(t)=2\int_0^\infty\Phi(u)\cos(tu)\,du$, so $\gamma(k)=2\,k!\,b_k/(2k)!$;
the Tur\'an inequalities below are the hyperbolicity of $J_{2,m-1}$). With the moments
$b_m=\int_0^\infty t^{2m}\Phi(t)\,dt$, Csordas, Norfolk and Varga~\cite{CNV} proved
the Tur\'an inequalities
\[
b_m^2>\frac{2m-1}{2m+1}\,b_{m-1}b_{m+1}\qquad(m\ge1),
\]
so the order-two defect
$\delta_2=-\bigl(b_m^2-\tfrac{2m-1}{2m+1}b_{m-1}b_{m+1}\bigr)$ is negative
unconditionally. The order-$k$ condition is the hyperbolicity of the degree-$k$
Jensen polynomials $J_{k,n}$; for each $k$ it is known for $n$ beyond an effective
shift $N(k)$~\cite{GORZ}. As said in Section~\ref{sec:line}, no bound on $N(k)$
uniform in $k$ is known.

\section{Defect-rate transfer}
The continuity lemma can be made quantitative.
\begin{proposition}[Rate transfer]\label{prop:rate}
Suppose $\|Q_m-Q_\infty\|\le\eta(m)$ and $\max_j\|\Pi_m^j-\Pi_\infty^j\|\le\eta(m)$
with $\eta(m)\to0$. Then
$|\delta_m-\delta_\infty|\le(1+4\|Q_\infty\|)\eta(m)+O(\eta(m)^2)$. In particular,
if $\delta_\infty>0$, then $\delta_m\ge\delta_\infty/2$ as soon as
$(1+4\|Q_\infty\|)\eta(m)+O(\eta(m)^2)<\delta_\infty/2$.
\end{proposition}
\begin{proof}
Write $q_m(v)=(-1)^jQ_m(v,v)$ on $V^j$. Let $w\in P_\infty^j$, $\|w\|=1$, and
$v=\Pi_m^jw\in P_m^j$. Then $\|v-w\|=\|(\Pi_m^j-\Pi_\infty^j)w\|\le\eta$, so
$1-\eta\le\|v\|\le1$. We have $|q_m(v)-q_\infty(v)|\le\eta\|v\|^2\le\eta$ and
$|q_\infty(v)-q_\infty(w)|\le\|Q_\infty\|(\|v\|+\|w\|)\|v-w\|\le2\|Q_\infty\|\eta$.
Moreover $|q_m(v)|\le\|Q_\infty\|+\eta$ and $\|v\|^{-2}\le1+2\eta+O(\eta^2)$, so
$q_m(v/\|v\|)\le q_m(v)+2\|Q_\infty\|\eta+O(\eta^2)$. Hence
$\delta_m\le q_m(v/\|v\|)\le q_\infty(w)+(1+4\|Q_\infty\|)\eta+O(\eta^2)$; taking
the infimum over $w$ and the minimum over $j$ gives
$\delta_m\le\delta_\infty+(1+4\|Q_\infty\|)\eta+O(\eta^2)$. For the reverse
inequality start from a unit $v\in P_m^j$ and set $w=\Pi_\infty^jv$; the same
estimates (with $\|Q_m\|\le\|Q_\infty\|+\eta$) give
$\delta_\infty\le\delta_m+(1+4\|Q_\infty\|)\eta+O(\eta^2)$.
\end{proof}

\section{Conclusion}
We separated eventual and uniform positivity (Theorem~\ref{thm:crit}): they
coincide for scale-free hierarchies, and a bounded threshold gives uniform
positivity beyond one scale. Non-trivial scale dependence with bounded threshold
exists (Proposition~\ref{prop:exist}), but convergence of measures does not in
general produce it (Example~\ref{ex:def}), nor a uniform gap
(Example~\ref{ex:deg}). For convergent
Hodge--Riemann towers the gap is continuous in the limit, with an explicit rate
(Lemma~\ref{thm:hr}, Proposition~\ref{prop:rate}). For $\xi$, no statement about
RH is proved; the known thresholds $N(k)$ are not known to be bounded uniformly
in $k$. A natural question is a uniform Hodge--Riemann
statement for families, such as matroids or polymatroids of growing size, to
which Lemma~\ref{thm:hr} does not apply.

\appendix
\section{Numerical illustration of the continuity lemma}\label{app:code}
The following program builds two convergent families of symmetric bilinear forms
on a fixed space --- one converging to a strictly definite limit, one to a
degenerate limit --- and tracks the gap $\delta_m$, illustrating
Lemma~\ref{thm:hr}. The families are constructed to have the stated gaps.

\lstinputlisting[language=Python,caption={\texttt{code.py}}]{code.py}

\medskip\noindent\textbf{Output.}
\begin{lstlisting}
 definite limit (delta_inf>0):  expect delta_m -> 0.5 (inf_m delta_m > 0)
   m=   2: delta=0.2010
   m=   8: delta=0.4120
   m=  32: delta=0.4767
   m= 128: delta=0.4941
   m= 512: delta=0.4985
 degenerate limit (delta_inf=0):  expect delta_m -> 0 (inf_m delta_m = 0; delta_m > 0 for every m)
   m=   2: delta=0.5000
   m=   8: delta=0.1250
   m=  32: delta=0.0312
   m= 128: delta=0.0078
   m= 512: delta=0.0020
\end{lstlisting}
In the definite family $\delta_m\to\delta_\infty=1/2$; the degenerate family is
$Q_\infty+\tfrac1mI$, so by construction $\delta_m=1/m$.

\section{A degenerate limit with positivity at every level}\label{app:deg}
For Example~\ref{ex:deg}, the following computes the smallest eigenvalue of the
order-$N$ Hankel matrix $[\int x^{a+b}d\nu_m]_{0\le a,b\le N}$ (size $N+1$) of $\nu_m=(1-\tfrac1m)\delta_0+\tfrac1m
\lambda_{[0,1]}\to\delta_0$.
\lstinputlisting[language=Python,caption={\texttt{code2.py}}]{code2.py}
\noindent\textbf{Output.}
\begin{lstlisting}
 order-N Hankel smallest eigenvalue of nu_m (degenerate limit nu_inf=delta_0):
        m    N=2   N=3   N=4
        4  1.86e-03  5.19e-05  1.51e-06
       16  4.89e-04  1.32e-05  3.80e-07
       64  1.24e-04  3.32e-06  9.53e-08
      256  3.10e-05  8.31e-07  2.38e-08
     1024  7.74e-06  2.08e-07  5.96e-09
 all eigenvalues positive; decay ~1/m at fixed N.
\end{lstlisting}
All eigenvalues are positive. At fixed order the smallest eigenvalue decreases
by a factor $4$ when $m$ is multiplied by $4$, i.e.\ it is $\asymp1/m$: the
positivity gap is not uniform in $m$, although positivity holds at every $m$.

\begin{thebibliography}{9}
\bibitem{AHK} K.~Adiprasito, J.~Huh, E.~Katz, \emph{Hodge theory for combinatorial geometries}, Ann. of Math. (2) \textbf{188} (2018), no.~2, 381--452.
\bibitem{BH} P.~Br\"and\'en, J.~Huh, \emph{Lorentzian polynomials}, Ann. of Math. (2) \textbf{192} (2020), no.~3, 821--891.
\bibitem{CNV} G.~Csordas, T.~S.~Norfolk, R.~S.~Varga, \emph{The Riemann hypothesis and the Tur\'an inequalities}, Trans. Amer. Math. Soc. \textbf{296} (1986), 521--541.
\bibitem{GORZ} M.~Griffin, K.~Ono, L.~Rolen, D.~Zagier, \emph{Jensen polynomials for the Riemann zeta function and other sequences}, Proc. Natl. Acad. Sci. USA \textbf{116} (2019), no.~23, 11103--11110.
\bibitem{RP06} D.~A.~Trejo Pizzo, \emph{Positivity formulations of the Riemann hypothesis and the function-field analogy: a heuristic checklist and some remarks}, preprint.
\bibitem{RP14} D.~A.~Trejo Pizzo, \emph{Hodge-index and hyperbolicity reformulations of the Riemann hypothesis as positivity statements}, preprint.
\end{thebibliography}

\end{document}
